\documentclass[12pt, a4paper]{article}
\usepackage[utf8]{inputenc}
\usepackage[T1]{fontenc}
\usepackage[english]{babel}
\usepackage{amsmath}
\usepackage{amssymb}
\usepackage{amsthm}
\usepackage{mathtools}
\usepackage[
  top=2cm,
  bottom=2cm,
  left=2cm,
  right=2cm
]{geometry}
\usepackage{setspace}
\onehalfspacing

% Graphics & floats
\usepackage{graphicx}
\usepackage{float}
\usepackage{caption}
\usepackage{subcaption}

\setlength{\parindent}{0pt}

% References & links
\usepackage[
  backend=biber,
  style=alphabetic,
  sorting=nyt
]{biblatex}
\addbibresource{references.bib}
\usepackage[
  colorlinks=true,
  linkcolor=black,
  citecolor=blue,
  urlcolor=blue
]{hyperref}

% Miscellaneous
\usepackage{microtype}
\usepackage{enumitem}
\usepackage{booktabs}

\theoremstyle{plain}
\newtheorem{theorem}{Theorem}[section]
\newtheorem{lemma}[theorem]{Lemma}
\newtheorem{proposition}[theorem]{Proposition}
\newtheorem{corollary}[theorem]{Corollary}

\theoremstyle{definition}
\newtheorem{definition}[theorem]{Definition}
\newtheorem{example}[theorem]{Example}
\newtheorem{remark}[theorem]{Remark}

\newcommand{\R}{\mathbb{R}}
\newcommand{\C}{\mathbb{C}}
\newcommand{\N}{\mathbb{N}}
%\newcommand{\i}{\text{i}}
\newcommand{\Eeps}{E_\varepsilon}
\newcommand{\psieps}{\psi_\varepsilon}
\newcommand{\ustar}{u^*}
\newcommand{\psistar}{\psi^*_\varepsilon}
\newcommand{\Wren}{W}
\newcommand{\Jac}{J\psi}
\newcommand{\sj}{j(\psi)}
\newcommand{\Lap}{\Delta}
\newcommand{\dd}{\,\mathrm{d}}
\newcommand{\norm}[1]{\left\lVert #1 \right\rVert}
\newcommand{\abs}[1]{\left\lvert #1 \right\rvert}
\newcommand{\inner}[2]{\left\langle #1,\, #2 \right\rangle}

% Shorthand for the symplectic matrix
\newcommand{\symp}{J_{\mathrm{symp}}}

\title{
  \Large\textbf{Numerical Simulation of Quantized Vortices\\
  with Small Core Size via Vortex Tracking}\\[0.5em]
  \large Complex Quantum Systems and Vortex Dynamics\\
  KIT --- Fakultät für Mathematik, Sommersemester 2026
}
\author{Thomas Ackermann}
\date{\today}

\begin{document}

\maketitle
\thispagestyle{empty}

\section{Introduction}

Ginzburg-Landau Hamiltonian with core size $\varepsilon$:
\begin{align}
    E_\varepsilon (u) = \int_\Omega \frac{1}{2} |\nabla u|^2 + \frac{1}{4 \varepsilon^2} (1- |u|^2)^2
\end{align}
\begin{itemize}
    \item GP PDE: $i \partial_t \psieps = \Delta \psieps + \frac{1}{\varepsilon^2} (1 - |\psieps|^2) \psieps$
    \item Ansatz for solution: $\phi_\varepsilon (x) = f_{\varepsilon, r_0} (|x|) \exp (i \theta(x))$
    \item $f_{\varepsilon, r_0}(0) = 0$, $f(\varepsilon, r_0) (r_0)= 1$
    \item draw graph
\end{itemize}
Renormalized energy
\begin{itemize}
    \item Ginzburg-Landau energy diverges as $\varepsilon \to 0$.
    \item This divergence comes purely from the self-energy - no information about vortex interactions
    \item After subtracting the divergence we get the renormalized energy
\end{itemize}

\begin{align}
  W(a,d) =  - \pi \sum_{1 \leq i \neq j \leq N} d_i d_j \log |a_i - a_j| - \pi \sum_j d_j R(a_j ; a, d)
\end{align}
\begin{itemize}
    \item first term: vortex-vortex interaction
    \item second term: vortex-boundary interaction
    \item $R$ is harmonic correction that enforces the Neumann boundary condition
\end{itemize}

Important Quantities:
\begin{itemize}
  \item Supercurrents $j(\psi) = \frac{1}{2i}(\overline{\psi} \nabla - \psi \nabla \overline{\psi})$
  \item Vorticity: $J \psi = \frac{1}{2} \nabla \times j(\psi)$
\end{itemize}

For $\varepsilon \to 0$:
$J \psieps (t) \to \pi \sum^N_{j=1} d_j \delta_{a_j (t)}$
The singular points solve:
\begin{align}
  &\dot a_j (t) = - \frac{1}{\pi} d_j J \nabla_{a_j} W(a(t), d) \\
  &a_j(0) = a^0_j
\end{align}

\begin{align}
    \nabla_{a_j} W(a,d) = - 2\pi d_j \nabla_x (R(x; a,d) + \sum^n_{j\neq k} d_k \log |x-a_k|)\vert_{x=a_j}
\end{align}
Boundary terms in the renormalized energy $R(x; a,d)$ which solves
\begin{align}
  &\Delta R = 0 \text{, in } \Omega \\
  &R = - \sum^N_{j=1} d_j \log |x-a_j| \text{, on } \partial \Omega
\end{align}

\begin{itemize}
    \item Following: condtions that determine the supercurrents
    \item They in turn determine harmonic map $u^*$
\end{itemize}
\begin{align}
  &\nabla j(u^*) =0\\
  &\nabla \times j(u^*) = 2\pi \sum^N_{j=1} d_j \delta_{a_j} \\
  &\nu j(u^*) =0 \text{ on } \partial \Omega
\end{align}
Where $\nu$ is the direction of the outward normal of $\partial \Omega$.\\ 
$u^*$ is then given by
\begin{align}
  u^* (x; a,d) = \exp (i H(x)) \prod^N_{j=1} \frac{x-a_j}{|x-a_j|}^{d_j}
\end{align}
$H(x)$ is a harmonic function that is determined such that $u^*$ satisfies the Neumann boundary condition.

\begin{definition}
  A familiy of inital conditions $(\psieps^0)_{\varepsilon>0}$ is said to be well-prepared if it satisfies the following assumptions,
  for some constant $C>=0$, $0<\alpha<1$ and $\varepsilon$ small enough:
  \begin{enumerate}
    \item There exists $N$ vortices with positions
    $a^0 = (a^0_j)_{j=1, \ldots , N} \in \Omega^{*N}$
    and degrees $d= (d_j)_{j=1, \ldots, N} \in \{\pm 1\}^N$
    such that
    \begin{align}
      ||J\psieps^0 - \pi \sum^N_{j=1} d_j \delta_{a_j^0}|| \leq C \varepsilon^\alpha 
    \end{align}
    and the vortices are distant enough
    \item the energy of $\psieps^0$ is close to be optimal:
    \begin{align}
      E_\varepsilon (\psieps^0) \leq W_\varepsilon (a_0, d) + C \varepsilon^{1/2}
    \end{align}
  \end{enumerate}
\end{definition}
$\dot W^{-1, 1}$ is the dual space of $W^{1, \infty}$
The norm is given by:
\begin{align}
    ||\mu||_{\dot W^{-1, 1}} = \sup \{\int_\Omega \varphi \dd \mu : ||\nabla \varphi||_{L^\infty} \leq 1, \varphi \in W^{1, \infty}_0 (\Omega)\}
\end{align}

\begin{theorem}
    Let $\psieps$ solve the equation corresponding to the Schrödinger flow of the Ginzburg energy, with well-prepared initial conditions
    for some initial vortices with positions $a^0$ and degrees $d_j$.
    Then there exists $\gamma < 1$ and $C > 0$, such that for any $\varepsilon < \varepsilon_0$,
    well-preparedness is preserved along time,
In particular,
\begin{align}
    &||J \psieps (t) - \pi \sum^N_{j=1} d_j \delta_{a_j (t)} ||_{\dot W^{-1, 1}} \leq C \varepsilon^\gamma \\
    & || j(\psieps (t)) - j(u^* (a(t), d))_{L^{4/3}} \leq C \varepsilon^\gamma
\end{align}
\end{theorem}

\begin{itemize}
    \item Question: How do we construct a well-prepared family?
\end{itemize}

For $\varepsilon$, there is a unique minimizer
$\phi_{\varepsilon, r_0} (x) = f_{\varepsilon, r_0} (|x|) \exp (i \theta(x))$
where $f_{\varepsilon, r^0}$ satisfies an ODE, which we have seen before. 

Starting from the localized approximation $\phi_\varepsilon$ of a single vortex one can the build the following inital condition:
\begin{align}
    \psieps^* (x; a^0, d) = u^* (x; a^0, d) \prod^N_{j=1} f_\varepsilon (|x- a^0_j|), \quad x \in \overline \Omega
\end{align}
which has the well preparedness properties.

\section{Method}
\textbf{Idea}
\begin{itemize}
    \item FEM requires fine mesh for small $\varepsilon$
    \item here $\varepsilon is not the limiting factor$
    \item we can use the point vortex system to track the vortex positions
    \item we can then build an approximation of the wave function from the vortex positions and use
    it as an initial condition for the PDE solver
\end{itemize}
$R$ can be solved numerically like this:
\begin{enumerate}
\item Choose a maximum degree $n$ and fix it
\item Compute the Fourier coefficients of the Dirichlet boundary condition up to order $n$ for instance using a FFT
\item Compute the approximate solution $R_n$ as the harmonic expansion of $\mathbb{P}_n g_a$
\begin{align}
    R_n (r \exp (i \theta)) = \sum^n_{k=-n} r^{|k|} \hat g_a (k) \exp (i k \theta)
\end{align}
\end{enumerate}

\subsection{Algorithm}

\subsubsection{Step 1: Constructing well-prepared initial conditions}
\begin{itemize}
\item Define initial vortex positions and degrees. Set up the initial phase $H$ such that the Neumann boundary condition is satisfied.
\item Compute the radial function $f_\varepsilon$
\end{itemize}
\subsubsection{Step 2: evolve vortex positions via Hamiltonian ODE}

solve the ODE using RK4
and compute $R$ as described before.
\\
\subsubsection{Computing $R$}
\begin{itemize}
    \item The solution $R$ is harmonic and the boundary conditions are smooth as long as the vortices stay away from the boundary.
    \item We use harmonic polynomials as a basis to solve the PDE.
    \item For simplicity we choose as the unit disk.
\end{itemize}
harmonic polynomials:
\begin{align}
    \begin{cases}
        h_1 = 1 \\
        h_{2m} = \Re (x+ i y)^m \\
        h_{2m + 1} =  \Im (x+ iy)^m
    \end{cases}
\end{align}
for $x,y \in \mathbb \R$.

\begin{itemize}
    \item When $\Omega$ is the unit disk, the restriction of these polynomials to the boundary is nothing else than the Fourier basis.
    \item Denote by $\mathbb{P}_n$ the orthogonal projection from $L^2 \partial(\Omega)$ to $\mathcal X_n := \text{span}\{h_k \mid_{\partial\Omega}, i \leq k \leq 2n+1\}$.
\end{itemize}

\begin{itemize}
    \item In polar coordinates: $\Delta R = \frac{1}{r} \partial_r (r \partial_r R) + \frac{1}{r^2} \partial_{\theta}^2 R$
    \item Ansatz: $R(r,\theta) = f(r) \exp (i k \theta)$ for some $k \in \mathbb Z$.
    \item This gives us the ODE: $f''(r) + \frac{1}{r} f'(r) - \frac{k^2}{r^2} f(r) = 0$
    \item Using the Ansatz $f(r)= r^\alpha$ you get $\alpha = \pm |k|$
    \item one gets: $R_k(r, \theta) = r^{|k|} \exp(i k \theta)$
    \item PDE is linear therefore linear combinations are solutions
    \item on the unit circle for harmonic functions one gets: $R(r, \theta) = \sum^\infty_{-\infty} c_k r^{|k|} \exp(i k \theta)$
    \item the coefficients are determined by the boundary conditions
    \item At the boundary: $R(1, \theta) = \sum^\infty_{-\infty} c_k \exp(i k \theta) = g(\theta)$
    \item this is just the Fourier expansion of $g$
    \item therefore $c_k = \frac{1}{2\pi} \int^{2\pi}_0 g(\theta) \exp (- i k \theta) \dd \theta$
\end{itemize}
\subsubsection{Step 3: Smoothing}
At time $t>0$ build back an approximation of the wave function from the vortex positions $a(t)$ as
\begin{align}
    u^* (x; a(t), d) = \exp (i H(x)) \prod^N_{j=1} (\frac{x-a_j}{|x-a_j|})^{d_j}
\end{align}
\begin{itemize}
    \item $u^*$ has singularities at the vortex locations $a_j$
    \item This is the reason we need to smooth it out before using it as an initial condition for the PDE
    \item Key insight: $\epsilon$ only enters through the 1D preprocessing - it is no longer a discretization bottleneck
    \item $H$ solves: $\Delta H = 0$ in $\Omega$ and $H = - \sum^N_{j=1} d_j \partial_\tau \log |x-a_j|$ on $\partial \Omega$
    \item $\tau$ is the direction of the tangent to $\partial \Omega$
    \item $H$ can also be solved using the same harmonic polynomial basis
\end{itemize}

\section{Error Analysis}

\begin{lemma}
  Let $\Omega \subset \R^2$ be the unit disk, let $\{a(t)\}_{t\geq 0}$ denote the exact solution to the Hamiltonian mechanics.
  Let $b(t)$ be the approximated ODE trajectory with the same degrees and initial positions $b^0$ with harmonic polynomials of degree up to $n$
  and time step $\delta_t$. Suppose that for some $T>0$,
  \begin{align}
  \rho_T = \min_{t\leq T} \min_{j\neq k} \{ |b_j (t) - b_k  (t)|, \text{dist} (b_k(t) , \partial \Omega), |a_j (t) - a_k  (t)|, \text{dist} (a_j(t) , \partial \Omega)\}>0
  \end{align}
  is such that $0 < \rho_T < 1$. Then we have:
  \begin{align}
    |a(t) - b(t)| \leq C (|a^0 -b^0| + (1- \rho_T)^n n^{1-l}) + C \delta t^4
  \end{align}
\end{lemma}

In order to prove this we need 2 other lemmas 

\begin{lemma}
  Let $b^0, a^0 \in \Omega^{*N}$, for fixed $d \in \{\pm1\}^N$
  and let $\rho(a) := \min_{j\neq l} \{\text{dist}(a_j, \partial\Omega) , |a_j - a_l|\}$.
  Suppose that $\rho_T:= \text{inf}_{t\in [0,T]} \leq 0$
  \begin{align}
  |a(t) - b(t)| \leq  C (|a^0-b^0| + \frac{\text{sup}_{s\in [0,T]}||\mathbb{P}^\perp_n g_{a(s)}||_{L^2(\partial \Omega)}}{\rho^{3/2}_T})\exp{C t /\rho^{5/2}_T}
  \end{align}
  Where $g_a(x) = - \sum^N_j d_j \log |x-a_j|$
\end{lemma} 
\begin{proof}
    \begin{itemize}
\item $|F^n(b) - F(a)| \leq E_1 + E_2 + E_3$
\item $E_1 = \sum_j^N |\nabla_x R^n(b_j; b ,d) - \nabla_x R^n(a_j; a, d)|$
\item $E_2 = \sum_j^N |\nabla_x R^n (a_j; a, d) - \nabla R (a_j; a, d)|$
\item $E_3 = \sum_{j \neq k} |\frac{b_j - b_k}{|b_j - b_k|^2} - \frac{a_j - a_k}{|a_j - a_k|^2}|$
\item $E_1$ is the error made takinig different evaluation points on the boundary, $E_2$ discretizatio error, $E_3$ is the error made by the vortex-vortex interaction term
    \end{itemize}

    \textbf{$E_3$ error bound}
    \begin{itemize}
        \item $E_3$ is the error made by the vortex-vortex interaction term
        \item $|\nabla \log |x-a| - \nabla \log |x-b|| \leq C \frac{|a-b|}{\min \{ |x-a|, |x-b|\}}^2$
        \item set $x = a_j$, $a=a_k$, $b=b_k$: $E_3 \leq C \frac{a-b||}{\min \{\rho(a), \rho(b)\}^2}$
        \item $\rho(a) = \min_{j\neq k} \{\text{dist}(a_j, \partial\Omega) , |a_j - a_k|\}$
    \end{itemize}
    \textbf{$E_2$ error bound}
    \begin{itemize}
        \item discretization error
        \item the difference $R^n - R$ solve the Laplace equation with boundary data $P^\perp_n g_a$
        \item $\mathbb P^\perp_n$ is the orthogonal complement of $\mathbb P_n$
        \item this is everything that the fourier expansion of $g_a$ up to order $n$ does not capture $P^\perp_n g_a = g_a - P_n g_a = g_a - \sum^n_{k=-n} \hat g_a(k) \exp (ik\theta)$
        \item $\Delta (R^n - R) = 0$ in $\Omega$, $R^n - R = P^\perp_n g_a$ on $\partial \Omega$
        \item One can show that $E_2(a) \leq \frac{||\mathbb P^\perp_n g_a||_{L^2(\partial \Omega)}}{\min_j \{\text{dist} (a_j , \partial \Omega)\}^\frac{3}{2}}$ 
    \end{itemize}
    \textbf{$E_1$ error bound}
    \begin{itemize}
        \item $\sum_j^N \nabla_x R_n (a_j; b, d)$ and using triangle inequality
        \item $E_1 (a,b) \leq \frac{|a-b|}{ \min_j \{\text{dist} (a_j , \partial \Omega), \text{dist}(b_j ,\partial \Omega)\}}$
    \end{itemize}
    Putting everything together we get:
    \begin{align}
        |\dot a(t) - \dot b(t)| \leq C (\frac{|a(t) - b(t)|}{\min \{\rho(a(t)), \rho(b(t))\}^2} + \frac{||\mathbb P^\perp_n g_{a(t)}||_{L^2(\partial \Omega)}}{\min_j \{\text{dist} (a_j , \partial \Omega)\}^\frac{3}{2}})
    \end{align}
    \begin{itemize}
        \item Using the assumption on $\rho_T$ we get:
        \begin{align}
            |\dot a(t) - \dot b(t)| \leq C (\frac{|a(t) - b(t)|}{\rho_T^2} + \frac{||\mathbb P^\perp_n g_{a(t)}||_{L^2(\partial \Omega)}}{\rho_T^{3/2}})
        \end{align}
        \item Inequality: $\frac{d}{dt} |a-b| \leq |\dot a - \dot b|$
        \item Proof of inequality: $\frac{d}{dt} |a-b| = \frac{d}{dt} \sqrt{(a-b) \cdot (a-b)} = \frac{(a-b) \cdot (\dot a - \dot b)}{|a-b|}$
        \item use Cauchy Schwarz: $(a-b)\cdot (\dot a - \dot b) \leq |a-b| \cdot |\dot a - \dot b|$
        \item Gronwall-inequality: if
        \begin{align}
            \frac{d}{dt} y(t) \leq f(t) y(t) + g(t),
        \end{align}
 then $y(t) \leq y(0) \exp (\int^t_0 f(s) ds) + \int^t_0 g(s) \exp (\int^t_s f(\tau) d\tau) ds$
        \item Using Gronwall's inequality we get the desired bound
        \end{itemize}
\end{proof}


The time integration errors follows from the previous bounds and standard considerations on the numerical integration of ODEs

\begin{lemma}
  Let $b$ be the exact solution to $\dot{b} = F^n (b)$ and let $(b^k_{\delta t})_{k \delta t \leq T}$ be its numerical approximation with RK4 method and time step $\delta t$
  \begin{align}
    b^0_{\delta t} = b^0\\
    b^0_{\delta t} = b^{k-1}_{\delta t} + \delta t \Phi_{\delta t , n} (b^{k-1}_{\delta t})
  \end{align}
Where $\Phi_{\delta t, n}$ is the RK4 increment function associated with the approximate forcing term $F^n$.
Assume that:
\begin{align}
  0 < \rho_T = \min (\inf_{0 \leq t \ leq T} \rho(b(t)), \min_{k \delta t \leq T} \rho (b^k_{\delta t}))
\end{align}
Then there exists constants $C_{\rho_T}>0$ and $\Lambda_{\rho_T}$, that depends on $\rho_T$ but not on $n, b$ or $\delta t$ such that
the following in time error holds:
\begin{align}
  \max_{k \delta t < T} |b(k \delta t) - b^k_{delta_t}| \leq \delta t^4 \frac{C_{\rho_T}}{\Lambda_{\rho_T}} (\exp (\Lambda_{\rho_T} T)- 1)
\end{align}
\end{lemma}


\section{Numerical Results}
\begin{itemize}
    \item Show error increases with bigger $\varepsilon$
    \item Runtime: few seconds on laptop for $\varepsilon = 0.01$ and $N=10$ vortices
    \item FEM solver took about a week to run (non-optimized code on a small cluster)
    \item Parameters of simulation: mesh size: $2 \cdot 10^{-3}$, time step: $\delta t = 10^{-5}$
\end{itemize}

\newpage
\printbibliography